\documentclass{article}

\usepackage{iclr2027_conference,times}
\usepackage[T1]{fontenc}
\usepackage{lmodern}
\usepackage{amsmath,amssymb,amsthm,mathtools}
\usepackage{booktabs,microtype,url,xcolor,graphicx,algorithm,algorithmic}
\usepackage[colorlinks=true,allcolors=blue]{hyperref}

\newtheorem{theorem}{Theorem}
\newtheorem{lemma}[theorem]{Lemma}
\newtheorem{proposition}[theorem]{Proposition}
\newtheorem{corollary}[theorem]{Corollary}
\theoremstyle{definition}
\newtheorem{definition}[theorem]{Definition}
\theoremstyle{remark}
\newtheorem{remark}[theorem]{Remark}

\newcommand{\E}{\mathbb E}
\newcommand{\R}{\mathbb R}
\newcommand{\DeltaK}{\Delta_K}
\newcommand{\Dir}{\operatorname{Dir}}
\newcommand{\TV}{\operatorname{TV}}
\newcommand{\Reg}{\operatorname{Reg}}
\newcommand{\PUCal}{\operatorname{PUCal}}
\newcommand{\supp}{\operatorname{supp}}

\title{Dirichlet Follow-the-Leader Closes the Gap in\\Simultaneous Multiclass U-Calibration}
\author{\textbf{Pahan Dewasurendra} \\ Johns Hopkins University \\ August 4, 2026}

\begin{document}
\maketitle

\begin{abstract}
Can one forecaster attain the optimal regret rate for every bounded proper loss
and also adapt to every smooth proper loss?  Recent work answered this up to a
dimension gap.  Its self-concordant perturbation gives roughly
$K^{5/4}\sqrt T$ worst-case regret and incurs an additional
$\beta\sqrt K\log K$ for $\beta$-smooth losses.  We close both gaps with a
one-line forecaster.  After observing class counts $c_{t-1}$, draw the next
prediction from $\Dir(c_{t-1})$, on the face of classes seen so far.  This is a
fresh Bayesian bootstrap of the outcomes.  The analysis rests on an exact
identity: averaging any bounded proper loss under $\Dir(\alpha)$ equals a
discrete derivative of its Dirichlet-averaged Bayes risk.  The identity makes
the be-the-perturbed-leader term telescope to a nonpositive Jensen gap.  A
one-count likelihood ratio then bounds stability by the inverse square root
of that class's count.  The resulting single, horizon-free algorithm satisfies
\begin{align*}
 \sup_{\ell}\E\Reg_\ell&\leq4\sqrt{S_TT}\leq4\sqrt{KT},\\
 \E\Reg_\ell&\leq\tfrac52\beta(1+\log T)
 \quad\text{for every $\beta$-smooth proper $\ell$.}
\end{align*}
Here $S_T$ is the number of observed classes.  Known lower bounds show that
both rates are optimal in their nontrivial regimes.  The proof covers
nondifferentiable losses and changes of the active simplex face.
\end{abstract}

\section{Introduction}

A probabilistic forecast is often consumed by an agent whose downstream
utility is unknown to the forecaster.  Proper losses capture this uncertainty:
each proper loss corresponds to a way of valuing probabilistic predictions,
and truthful prediction minimizes expected loss.  U-calibration asks for one
online sequence of forecasts with low regret under every bounded proper loss
\citep{kleinberg2023ucal}.  This requirement is substantially stronger than
optimizing a fixed score such as Brier loss.

For $K$ outcomes, \citet{luo2024optimal} established the optimal
$\Theta(\sqrt{KT})$ worst-case pseudo-U-calibration rate.  Their optimal
forecaster does not adapt to easier scores.  In particular, it can incur
$\Omega(\sqrt T)$ regret on squared loss even though Follow-the-Leader (FTL)
has logarithmic regret.  Very recent work showed that this conflict is not
fundamental \citep{frongillo2026simultaneous}.  A carefully shaped
self-concordant perturbation simultaneously gives logarithmic regret on
smooth scores and sublinear regret on all proper scores.  Its bounds leave two
linked gaps.  The general rate is $\widetilde O(K^{5/4}\sqrt T)$ rather than
$O(\sqrt{KT})$, and the smooth rate contains an additive
$O(\beta\sqrt K\log K)$ term.  The paper explicitly asks whether simultaneous
optimality is possible.

We answer this question affirmatively.  Let $c_t=\sum_{s\leq t}y_s$ be the
vector of outcome counts.  At time $t\geq2$, independently sample
\[
 P_t\sim\Dir(c_{t-1}).
\]
Zero count coordinates are fixed at zero, so the distribution lives on the
face spanned by classes observed so far.  Equivalently, assign independent
Gamma weights to the observed outcomes and normalize.  Thus the algorithm is
exactly a fresh Bayesian bootstrap \citep{rubin1981bootstrap}.  It requires no
horizon, smoothness parameter, learning rate, or optimization oracle.
The Bayesian-bootstrap draw itself is classical.  Our claims concern its new
adversarial proper-loss analysis and simultaneous regret guarantees.

The main proof is not a generic posterior-sampling argument.  If $f$ is the
concave Bayes risk of a proper loss and
$F(\alpha)=\E_{P\sim\Dir(\alpha)}f(P)$, we prove
\begin{equation}
 \E_{P\sim\Dir(\alpha)}\ell(P,e_j)
 =nF(\alpha)-(n-1)F(\alpha-e_j),
 \qquad n=\sum_i\alpha_i.                 \tag{1}\label{eq:key-intro}
\end{equation}
This identity applies to arbitrary bounded proper losses, including
polyhedral and V-shaped losses.  It turns the perturbed-leader contribution
into $T(F(c_T)-f(c_T/T))\leq0$.  The remaining stability term compares two
Dirichlet laws whose parameters differ by one count.  Their likelihood ratio
is linear in one coordinate, which gives stability $O(1/\sqrt m)$ when that
coordinate has appeared $m$ times.  Summing separately within each class
produces $O(\sum_i\sqrt{N_i})=O(\sqrt{KT})$.

The same distribution is automatically adapted to smooth losses.  Its mean
is empirical FTL and its squared radius is at most $1/t$.  Centering removes
the first-order smoothness term, leaving $O(\beta/t)$.  A short
semiconcavity argument also gives an explicit $2\beta H_T$ FTL bound under
the one-sided smoothness definition used in prior work.

Our contributions are:
\begin{itemize}
 \item We give one efficient, horizon-free forecaster with
 $4\sqrt{S_TT}\leq4\sqrt{KT}$ expected regret for every bounded proper loss and simultaneous
 $O(\beta\log T)$ regret for every bounded $\beta$-smooth proper loss.
 \item We prove the count-deletion identity in \eqref{eq:key-intro}, including
 nondifferentiable Bayes risks, zero count coordinates, and the singular case
 $\alpha_j=1$.
 \item We isolate two exact geometric facts behind simultaneous adaptation:
 one-count Dirichlet total variation controls nonsmooth loss, while centered
 covariance controls smooth loss.
\end{itemize}

We emphasize the quantifier.  Our result controls
$\sup_\ell\E\Reg_\ell$, called pseudo-U-calibration.  It does not claim the
stronger $\E\sup_\ell\Reg_\ell$ for the infinite class of all proper losses.
This is the same expected-regret criterion studied by
\citet{frongillo2026simultaneous}.

\section{Setting and algorithm}\label{sec:setting}

Let $e_1,\ldots,e_K$ denote the vertices of $\DeltaK$.  At round $t$, a
forecaster chooses a possibly random $P_t\in\DeltaK$, then observes an outcome
$y_t\in\{e_1,\ldots,e_K\}$.  We first state results for an oblivious outcome
sequence.  Appendix~\ref{app:adaptive} gives the standard extension to a
nonanticipating adaptive adversary when fresh randomness is used each round.

A loss $\ell:\DeltaK\times\{e_i\}_{i=1}^K\to[-1,1]$ is \emph{proper} if
\[
 p\in\arg\min_{q\in\DeltaK}\E_{Y\sim p}\ell(q,Y)
 \quad\text{for every }p\in\DeltaK.
\]
For an outcome sequence, propriety makes its empirical distribution
$q_T=T^{-1}\sum_ty_t$ a hindsight minimizer.  Define
\[
 \Reg_\ell(T)=\E\sum_{t=1}^T\ell(P_t,y_t)-T f(q_T),
 \qquad f(p)=\sum_i p_i\ell(p,e_i).
\]
The expectation is over the forecaster.  Let $\mathcal L$ be all proper
losses with range in $[-1,1]$ and let
$\PUCal_T=\sup_{\ell\in\mathcal L}\Reg_\ell(T)$.

Following \citet{frongillo2026simultaneous}, a differentiable loss is
$\beta$-smooth when
\begin{equation}
 \ell(q,y)-\ell(p,y)
 \leq\langle\nabla_p\ell(p,y),q-p\rangle
       +\frac\beta2\|q-p\|_2^2                 \tag{2}\label{eq:smooth}
\end{equation}
for every $p,q\in\DeltaK$ and outcome $y$.
Gradients and differentiability on the closed simplex are understood relative
to its affine hull, including one-sided differentiability on boundary faces.

For a nonnegative parameter vector $\alpha$, write $A=\{i:\alpha_i>0\}$.
We define $\Dir(\alpha)$ on the face $\Delta_A$ by drawing independent
$G_i\sim\mathrm{Gamma}(\alpha_i,1)$ for $i\in A$, setting
$P_i=G_i/\sum_{k\in A}G_k$, and setting $P_i=0$ outside $A$.

\begin{algorithm}[t]
\caption{Dirichlet Follow-the-Leader}
\label{alg:dftl}
\begin{algorithmic}[1]
\STATE Initialize $c_0=0$ and predict $P_1=(1/K,\ldots,1/K)$.
\FOR{$t=1,2,\ldots,T$}
  \IF{$t\geq2$}
    \STATE Draw $P_t\sim\Dir(c_{t-1})$ independently of past draws.
  \ENDIF
  \STATE Observe $y_t$ and set $c_t=c_{t-1}+y_t$.
\ENDFOR
\end{algorithmic}
\end{algorithm}

\begin{theorem}[Simultaneously optimal regret]\label{thm:main}
For every $K,T\geq1$, let $S_T$ be the number of distinct outcomes observed.
Algorithm~\ref{alg:dftl} satisfies
\[
 \PUCal_T\leq4\sqrt{S_TT}\leq4\sqrt{\min\{K,T\}T}.
\]
Simultaneously, for every $\beta$-smooth $\ell\in\mathcal L$,
\[
 \Reg_\ell(T)\leq\frac52\beta H_T
 \leq\frac52\beta(1+\log T),
\]
where $H_T=\sum_{t=1}^T1/t$.
\end{theorem}

For $K\leq T$, the $\sqrt{KT}$ worst-case rate is matched by the multiclass
V-shaped lower bound of \citet{luo2024optimal}.  Squared loss gives an
$\Omega(\log T)$ lower bound for the smooth subclass.  Hence both worst-case
rates in Theorem~\ref{thm:main} are optimal up to universal constants in their
nontrivial regimes.  The support-sensitive refinement is an additional
instance-dependent guarantee.

\section{A Dirichlet identity for every proper loss}\label{sec:identity}

Write $\ell_p=(\ell(p,e_1),\ldots,\ell(p,e_K))$.  Propriety gives
\begin{equation}
 f(q)\leq\langle q,\ell_p\rangle,
 \qquad f(p)=\langle p,\ell_p\rangle.          \tag{3}\label{eq:supergrad}
\end{equation}
Thus $f$ is concave and $\ell_p$ is a supergradient representation.  Since
$\|\ell_p\|_\infty\leq1$, applying \eqref{eq:supergrad} in both directions
also gives
\begin{equation}
 |f(p)-f(q)|\leq\|p-q\|_1.                    \tag{4}\label{eq:lipschitz-f}
\end{equation}

\begin{lemma}[Count deletion]\label{lem:identity}
Let $\alpha\geq0$, let $n=\sum_i\alpha_i$, and suppose $\alpha_j\geq1$.
For $F(\alpha)=\E_{P\sim\Dir(\alpha)}f(P)$,
\[
 \E_{P\sim\Dir(\alpha)}\ell(P,e_j)
 =nF(\alpha)-(n-1)F(\alpha-e_j).
\]
The statement uses the active-face convention above.  When $n=1$, the
second term is zero.
\end{lemma}

\begin{proof}
First suppose $\alpha_j>1$ and work relative to the active face.  A concave
function is differentiable almost everywhere on the relative interior.
At each such $p$, \eqref{eq:supergrad} makes the tangent component of
$\ell_p$ the unique supergradient of $f$.  Therefore
\begin{equation}
 \ell(p,e_j)=f(p)+D_{e_j-p}f(p)                \tag{5}\label{eq:savage-dir}
\end{equation}
almost everywhere.

Let $\rho_\alpha$ be the Dirichlet density and set $v(p)=e_j-p$.  A direct
calculation on the active face gives
\[
 \operatorname{div}v+D_v\log\rho_\alpha
 =\frac{\alpha_j-1}{p_j}-(n-1).
\]
The boundary flux vanishes.  At a face $p_i=0$ with $i\neq j$, the normal
component $v_i=-p_i$ cancels the possible density singularity.  At $p_j=0$,
the assumption $\alpha_j>1$ makes the flux vanish.  Integration by parts and
\eqref{eq:savage-dir} now yield
\begin{align*}
 \E_\alpha D_vf(P)
 &=(n-1)F(\alpha)
   -(\alpha_j-1)\E_\alpha\frac{f(P)}{P_j},\\
 (\alpha_j-1)\E_\alpha\frac{f(P)}{P_j}
 &=(n-1)F(\alpha-e_j),
\end{align*}
where the second equality is the ratio of Dirichlet normalizers.  Combining
the two displays proves the claim for $\alpha_j>1$.

If $\alpha_j=1$, apply the proved identity to
$\alpha+\varepsilon e_j$ and let $\varepsilon\downarrow0$.  The first
Dirichlet law converges in total variation, which is enough for the possibly
discontinuous bounded function $\ell(\cdot,e_j)$.  After deleting one count,
the law with parameter $\varepsilon$ in coordinate $j$ converges weakly to
the lower-dimensional face.  Equation~\eqref{eq:lipschitz-f} makes $f$
continuous, so its expectations converge.  The case $n=1$ is immediate from
$\ell(e_j,e_j)=f(e_j)$.
\end{proof}

The distinction between total variation and weak convergence in the last
paragraph matters.  A bounded proper loss can jump across a decision
boundary, while its Bayes risk remains continuous.

\section{Worst-case analysis}\label{sec:worst}

For analysis, after observing round $t$, draw a ghost prediction
$Q_{t+1}\sim\Dir(c_t)$.  For any fixed loss,
\begin{align}
 \Reg_\ell(T)
 &=\sum_{t=1}^T\E[\ell(P_t,y_t)-\ell(Q_{t+1},y_t)] \notag\\
 &\quad+\sum_{t=1}^T\E\ell(Q_{t+1},y_t)-Tf(q_T).  \tag{6}\label{eq:decomp}
\end{align}
If $y_t=e_j$, Lemma~\ref{lem:identity} gives
\[
 \E\ell(Q_{t+1},y_t)=tF(c_t)-(t-1)F(c_{t-1}).
\]
The second line of \eqref{eq:decomp} therefore telescopes to
\[
 T\{F(c_T)-f(q_T)\}\leq0,                     \tag{7}\label{eq:jensen}
\]
because $\E Q_{T+1}=q_T$ and $f$ is concave.

It remains to control stability.  The next elementary estimate is where the
coordinate adaptivity enters.

\begin{lemma}[One-count stability]\label{lem:tv}
Let $c\geq0$, $n=\sum_i c_i$, and $c_j=m>0$.  Then
\[
 \TV(\Dir(c),\Dir(c+e_j))
 \leq\frac12\sqrt{\frac{n-m}{m(n+1)}}
 \leq\frac1{2\sqrt m}.
\]
\end{lemma}

\begin{proof}
On their common active face, the likelihood ratio is
\[
 \frac{d\Dir(c+e_j)}{d\Dir(c)}(p)=\frac{np_j}{m}.
\]
Consequently,
\[
 \TV(\Dir(c),\Dir(c+e_j))
 =\frac12\E_{P\sim\Dir(c)}\left|\frac{nP_j}{m}-1\right|.
\]
Cauchy--Schwarz and
$\operatorname{Var}(P_j)=m(n-m)/(n^2(n+1))$ prove both inequalities.
When $b=n-m>0$, direct integration at the likelihood-ratio crossing
$p_j=m/n$ also gives the exact expression
\[
 \TV=\frac{(m/n)^m(1-m/n)^b}{mB(m,b)}.
\]
For $b=0$, both laws are the same point mass and the distance is zero.
\end{proof}

If class $j$ has appeared $m>0$ times before round $t$, the range
$[-1,1]$ and Lemma~\ref{lem:tv} bound the corresponding stability term by
$2\TV\leq1/\sqrt m$.  If it is the first appearance, the supports differ and
the trivial bound is $2$.  Let $N_i=c_{T,i}$ and
$S=|\{i:N_i>0\}|$.  Equations \eqref{eq:decomp} and \eqref{eq:jensen} give
\begin{align*}
 \Reg_\ell(T)
 &\leq2S+\sum_{i=1}^K\sum_{m=1}^{N_i-1}\frac1{\sqrt m}\\
 &\leq2S+2\sum_{i:N_i>0}\sqrt{N_i}\\
 &\leq2S+2\sqrt{ST}\leq4\sqrt{ST},
\end{align*}
where $S\leq T$ implies $S\leq\sqrt{ST}$.
The bound is independent of $\ell$, so taking the supremum proves the first
part of Theorem~\ref{thm:main}.

\section{Smooth-loss adaptation}\label{sec:smooth}

For $t\geq2$, Dirichlet moments give
\begin{equation}
 \E P_t=q_{t-1},\qquad
 \E\|P_t-q_{t-1}\|_2^2
 =\frac{1-\|q_{t-1}\|_2^2}{t}\leq\frac1t.     \tag{8}\label{eq:moments}
\end{equation}
Set $q_0=P_1$ for convenience.  Applying \eqref{eq:smooth}, taking
expectations, and using centering yields
\begin{equation}
 \E[\ell(P_t,y_t)-\ell(q_{t-1},y_t)]
 \leq\frac{\beta}{2t},\qquad t\geq2.           \tag{9}\label{eq:perturb-smooth}
\end{equation}

For completeness, we prove the FTL estimate under precisely the one-sided
definition \eqref{eq:smooth}.  This avoids assuming that \eqref{eq:smooth}
implies full gradient Lipschitzness.

\begin{lemma}[Smooth proper-loss FTL]\label{lem:smooth-ftl}
For every differentiable proper loss satisfying \eqref{eq:smooth}, empirical
FTL satisfies
\[
 \sum_{t=1}^T\ell(q_{t-1},y_t)-Tf(q_T)\leq2\beta H_T.
\]
\end{lemma}

\begin{proof}
We first establish an endpoint condition from propriety and boundary
differentiability.  Fix $i\neq j$ and set
$q_\varepsilon=(1-\varepsilon)e_j+\varepsilon e_i$.  Properness at
$q_\varepsilon$, comparing the truthful prediction with $e_j$, says
\[
 (1-\varepsilon)\ell(q_\varepsilon,e_j)
 +\varepsilon\ell(q_\varepsilon,e_i)
 \leq(1-\varepsilon)\ell(e_j,e_j)
 +\varepsilon\ell(e_j,e_i).
\]
Expanding at zero gives
$D_{e_i-e_j}\ell(e_j,e_j)\leq0$.  Propriety at $e_j$ also says that $e_j$
globally minimizes $\ell(\cdot,e_j)$, so the reverse inequality holds.
Thus every tangent derivative of $\ell(\cdot,e_j)$ at $e_j$ is zero.

Fix $p\in\DeltaK$, write $d=\|e_j-p\|_2$, and define
$\phi(s)=\ell((1-s)p+se_j,e_j)$.  Equation~\eqref{eq:smooth} is equivalent
to concavity of $h(s)=\phi(s)-(\beta d^2/2)s^2$.  Since $\phi'(1)=0$,
monotonicity of the derivative of $h$ gives, for $s<1$,
\[
 \phi'(s)-\beta d^2s=h'(s)\geq h'(1)=-\beta d^2.
\]
Hence $\phi'(s)\geq-\beta d^2(1-s)$.  The FTL update after outcome $e_j$ is
$q_t=(1-1/t)q_{t-1}+e_j/t$.  Integrating from $0$ to $1/t$ gives
\[
 \ell(q_{t-1},e_j)-\ell(q_t,e_j)
 \leq\frac{\beta\|e_j-q_{t-1}\|_2^2}{t}
 \leq\frac{2\beta}{t}.
\]
The standard be-the-leader inequality
$\sum_t\ell(q_t,y_t)\leq Tf(q_T)$ and summation prove the result.
\end{proof}

Adding \eqref{eq:perturb-smooth} to Lemma~\ref{lem:smooth-ftl} gives at most
$2\beta H_T+(\beta/2)(H_T-1)\leq(5\beta/2)H_T$, completing
Theorem~\ref{thm:main}.

\section{Related work and interpretation}\label{sec:related}

\paragraph{U-calibration.}
\citet{kleinberg2023ucal} introduced U-calibration for binary outcomes and
connected it to unknown downstream agents.  \citet{luo2024optimal} obtained
the optimal $\Theta(\sqrt{KT})$ multiclass pseudo-U-calibration rate, proved
the V-shaped lower bound, and identified loss classes on which FTL is
logarithmic.  Their worst-case-optimal perturbed leader is not adaptive to
smooth losses.  \citet{frongillo2026simultaneous} introduced a
self-concordant prediction-space perturbation that adapts to smooth and
log-barrier-smooth losses.  For the classes considered here, its bounds are
approximately $K^{5/4}\sqrt T$ and
$\beta\log T+\beta\sqrt K\log K$.  Our result removes these dimension gaps.
Calibeating, omniprediction, and swap-regret results study related universal
downstream guarantees with different benchmarks
\citep{roth2024forecasting,chen2026calibeating}.

\paragraph{Perturbation and posterior sampling.}
Follow-the-Perturbed-Leader is classical
\citep{kalai2005efficient,cesabianchi2006prediction}.  Fresh perturbations
also yield the standard nonanticipating adaptive-adversary extension
\citep{hutter2005adaptive}.  Algorithm~\ref{alg:dftl} can instead be viewed as
posterior sampling for a categorical model under the limiting zero-mass
Dirichlet prior, or as Rubin's Bayesian bootstrap.  Dirichlet-weighted loss
minimization appears in the loss-likelihood bootstrap
\citep{lyddon2019general}, and normalized random weights have also been
maintained in online Bayesian bagging \citep{lee2004lossless}.  Dirichlet Bayes
mixtures are classical in universal coding under log loss
\citep{watanabe2013achievability}, while size-biased Dirichlet integral
identities are classical probability machinery \citep{last2020integral}.
None of these works gives adversarial regret for an unknown proper loss.  Our
contribution is the loss-uniform count-deletion specialization, its
telescoping use, and the simultaneously optimal guarantees, not the bootstrap
draw itself.

\paragraph{Why Dirichlet geometry fits both regimes.}
The update $c\mapsto c+e_j$ has likelihood ratio $np_j/c_j$, so its total
variation automatically scales with the number of previous appearances of
the updated class.  This is the right geometry for discontinuous proper
losses.  At the same time,
\[
 \operatorname{Cov}(P_t)
 =\frac{\operatorname{diag}(q_{t-1})-q_{t-1}q_{t-1}^{\top}}{t},
\]
which has the centered $O(1/t)$ radius needed for smooth losses.  The exact
identity in Lemma~\ref{lem:identity} is what lets these two local facts control
global regret without a separate perturbed-leader path-length penalty.

\section{Limitations and conclusion}

Our theorem resolves simultaneous optimality for bounded proper losses and
the smooth subclass under the expected-regret criterion.  It does not control
actual U-calibration $\E\sup_{\ell\in\mathcal L}\Reg_\ell$ over the entire
infinite class.  Moving the supremum inside expectation requires complexity
control or a different argument.  We also do not address the broader class of
losses smooth relative to a log barrier, where Euclidean Dirichlet covariance
need not be the right measure.

The result shows that the previous tradeoff was geometric rather than
information-theoretic.  A Bayesian-bootstrap sample of empirical FTL has
exactly the count stability needed by nonsmooth scores and exactly the
centering needed by smooth scores.  The count-deletion identity makes those
two properties compatible in one short analysis.

\section*{AI use statement}
Generative AI tools were used to assist with literature search, symbolic and
numerical proof checks, and language editing.  The authors independently
checked the mathematical arguments, citations, and final manuscript and take
full responsibility for its content.

\bibliography{references}
\bibliographystyle{iclr2027_conference}

\appendix

\section{Adaptive outcomes}\label{app:adaptive}

The main text assumes a fixed outcome sequence.  The same expected bounds
hold for a nonanticipating adaptive adversary that may depend on past
predictions but not on the fresh draw $P_t$.  Condition on the history before
the round and on the selected outcome.  Given the current count vector,
$P_t$ is a fresh Dirichlet draw, so the one-step identity and stability bounds
hold conditionally.  Summing conditional expectations gives the same
telescoping count potential along the realized count path.  This is the usual
fresh-randomization reduction for perturbed leader
\citep{hutter2005adaptive}.  No guarantee is possible against an adversary
that observes the current randomized forecast before choosing the same-round
outcome under this protocol.

\section{Additional boundary details for Lemma~\ref{lem:identity}}

We record a standard approximation that makes the nonsmooth integration by
parts fully formal.  Restrict $f$ to the relative interior of the active face.
It is concave and Lipschitz by \eqref{eq:lipschitz-f}.  Convolve it on compact
subsets with a smooth mollifier, apply integration by parts to the smooth
approximants, and then exhaust the face.  The directional derivatives of a
concave Lipschitz function converge almost everywhere and are uniformly
bounded in tangent directions.  Dominated convergence applies to every term
when $\alpha_j>1$.  The boundary flux estimate in the main proof is uniform
under this exhaustion.

For $\alpha_j=1$, write
$\alpha^{(\varepsilon)}=\alpha+\varepsilon e_j$.  On the common active face,
the densities of $\Dir(\alpha^{(\varepsilon)})$ converge in $L^1$ to that of
$\Dir(\alpha)$, hence in total variation.  After deleting $e_j$, the marginal
coordinate $P_j$ has a Beta distribution with first parameter $\varepsilon$
and therefore converges to zero in probability.  Conditional proportions of
the remaining coordinates have distribution $\Dir(\alpha_{-j})$.  This proves
weak convergence to the deleted face.  Boundedness handles the left side and
continuity of $f$ handles the right side.

\section{Sanity checks}

An anonymous supplementary script evaluates the identity for Brier,
spherical, binary V-shaped, and multiclass polyhedral proper losses.  It also
checks the likelihood-ratio total-variation bound and exhaustively enumerates
all binary outcome sequences through $T=8$.  The smooth Brier identities are
also checked in closed form.  These calculations are not used by any proof.

\end{document}
